\section{Related Work}
\label{sec:related_work}

\textbf{Generalized semantic segmentation.}
The majority of existing semantic segmentation models are restricted to a fixed label set that is defined by the labels that are present in the training dataset~\citep{minaee2021image}. 
Few-shot semantic segmentation methods aim to relax the restriction of a fixed label set when one or a few annotated examples of novel classes are available at test time. These approaches learn to find reliable visual correspondences between a query image that is to be labeled and labeled support images that may contain novel semantic classes~\citep{shaban2017one,rakelly2018conditional,siam2019amp,wang2019panet,zhang2019pyramid,nguyen2019feature,liu2020part,wang2020few,tian2020prior,wang2020few, tian2020prior,boudiaf2021few, min2021hypercorrelation}. While this strategy can significantly enhance the generality of the resulting model, it requires the availability of at least one labeled example image with the target label set, something that is not always practical.

Zero-shot semantic segmentation approaches aim to segment unseen objects without any additional samples of novel classes. Text embeddings of class labels play a central role in these works. \citet{bucher2019zero} and \citet{Gu2020CaGNet} propose to leverage word embeddings together with a generative model to generate visual features of unseen categories, while \citet{xian2019semantic} propose to project visual features into a simple word embedding space and to correlate the resulting embeddings to assign a label to a pixel. \citet{hu2020uncertainty} propose to use uncertainty-aware learning to better handle noisy labels of seen classes, while \citet{li2020consistent} introduce a structured learning approach to better exploit the relations between seen and unseen categories. While all of these leverage text embeddings, our paper is, to the best of our knowledge, the first to show that it is possible to synthesize zero-shot semantic segmentation models that perform on par with fixed-label and few-shot semantic segmentation methods.

A variety of solutions have been proposed~\citep{zhang2020hybrid,liu2020few,perera2020generative,zhou2021learning} for open-set recognition~\citep{scheirer2012toward,geng2020recent}. These aim to provide a binary decision about whether or not a given sample falls outside the training distribution, but do not aim to predict the labels of entirely new classes.

Finally, a different line of work explores cross-domain adaptation methods for semantic segmentation by using feature alignment, self-training, and information propagation strategies~\citep{yang2021exploring,wang2021give}. The target of these works is to enhance the transferability of models to novel visual domains, but they do not address the issue of a restricted label set. As such they are orthogonal to our work.

\textbf{Language-driven recognition.}
Language-driven recognition is an active area of research. Common tasks in this space 
include visual question answering \citep{antol2015}, image captioning \citep{vinyals2014}, and image-text retrieval \citep{li20unienc}. CLIP~\citep{radford2021learning} demonstrated that classic recognition tasks that are not commonly associated with 
language can strongly benefit from language assistance. CLIP uses contrastive learning  together with high-capacity language models and visual feature encoders to synthesize extremely robust models for zero-shot image classification. Recent works have extended this basic paradigm to perform flexible object detection. ViLD~\citep{gu2021zero} introduces an advanced zero-shot object detection method that leverages CLIP, whereas MDETR~\citep{kamath2021mdetr} proposes an end-to-end approach that modulates a transformer-based baseline detector with text features that are obtained from a state-of-the-art language model. Like CLIP, these works have shown that the robustness and generality of object detection models can be strongly improved by language assistance. Our work is inspired by these approaches and presents, to the best of our knowledge, the first approach to flexibly synthesize zero-shot semantic segmentation models by leveraging high-capacity language models.

